\documentclass[main.tex]{subfiles}


\begin{document}

% \AddToShipoutPictureFG{%
% \begin{tikzpicture}[overlay,remember picture]
% \path (current page.south west) -- (current page.north east)
% node[midway,scale=1,color=gray,sloped,opacity=0.4] {\textnormal{Кравченко Игорь Игоревич\hspace{1cm} +79010144910\hspace{1cm} super.antig@yandex.ru}};
% \end{tikzpicture}
% }


% %from pagenumber to up
% \phantomsection 
% \hypertarget{toc}{}

 

 

% номер секции вручную
\setcounter{section}{20
}
\addtocounter{section}{-1} 

\section{Сила упругости}

\textbf{Деформация}~--- это изменение формы или размеров тела. На рис.\,\ref{fig:p38} изображены деформированные резиновый шар и пластилиновый брусок.

\begin{figure}[H]
    \centering

    \begin{tikzpicture}[font=\small,            line cap=round,                             >={Stealth[inset=0pt,round,length=3mm, angle'=20]},vec/.style={>={Stealth[inset=0pt,round,length=3.5mm, angle'=20]}}]


\filldraw[lightgray,semitransparent] (0,0) rectangle (5,-0.3);
\draw (0,0) -- (5,0);

% \filldraw[lightgray,draw=black] ({5/3-.5},0) rectangle++ (1,.5);

% \shade[ball color=lightgray] ({5/3},0.3) circle (.3);
%%%%%%%%%%%%%%%%%%%%%%%%%%%%%%%%%%%%%%

\filldraw[lightgray,draw=black] ({2*5/3-.5},0) --++ (1,0) --++ (0,.5) to [out=180,in=0] ++ (-.5,-.15) node (a) {} to [out=180,in=0] ++ (-.5,.15) --++ (0,-.5);

% \path ({5/3},0) arc (-90:155:.3cm) node (s) {};
% \shade[ball color=lightgray] ({2*5/3},0) arc (-90:25:.3cm) to [out=115,in=0] ({2*5/3},.35) node (b) {} to [out=180,in=65] (s.center) arc (155:270:.3cm);
\path ({5/3-.15},0) to [out=175,in=-110]++ (-.3,.3) node (s1) {};
\shade[ball color=lightgray] ({5/3+.15},0) to [out=5,in=-70]++ (.3,.3) to [out=110,in=0] ({5/3},.35) node (b) {} to [out=180,in=70] (s1.center) to [out=250,in=175] ({5/3-.15},0);
% \draw[] ({5/3},0.3) circle (.3);

%vec
\draw[->,red,thick,vec] (a.center) --++ (0,-1) node[black,right] {$\vec F$};
\node at (a.center) [circle,fill=red,inner sep=0pt,minimum size=1mm,label=above:Б] {};

\draw[->,red,thick,vec] (b.center) --++ (0,-1) node[black,right] {$\vec F$};
\node at (b.center) [circle,fill=red,inner sep=0pt,minimum size=1mm,label=above:Ш] {};

% \draw (b) --++ (-3,0);


    \end{tikzpicture}
\caption{Деформированные тела}
\label{fig:p38}
\end{figure}

Резиновый шар~Ш и пластилиновый брусок~Б испытывают деформации под действием некоторой внешней силы~$\vec F$, приложенной к телам сверху. 

Теперь внешние силы, деформирующие тела, сняли (рис.\,\ref{fig:p39}).

\begin{figure}[H]
    \centering

    \begin{tikzpicture}[font=\small,            line cap=round,                             >={Stealth[inset=0pt,round,length=3mm, angle'=20]},vec/.style={>={Stealth[inset=0pt,round,length=3.5mm, angle'=20]}}]


\filldraw[lightgray,semitransparent] (0,0) rectangle (5,-0.3);
\draw (0,0) -- (5,0);

% \filldraw[lightgray,draw=black] ({5/3-.5},0) rectangle++ (1,.5);

\shade[ball color=lightgray] ({5/3},0.3) circle (.3);
%%%%%%%%%%%%%%%%%%%%%%%%%%%%%%%%%%%%%%

\filldraw[lightgray,draw=black] ({2*5/3-.5},0) --++ (1,0) --++ (0,.5) to [out=180,in=0] ++ (-.5,-.15) node (a) {} to [out=180,in=0] ++ (-.5,.15) --++ (0,-.5);

%descr
\node[above] at (a) {Б};
\node[above] at ({5/3},0.6) {Ш};


% \path ({2*5/3},0) arc (-90:155:.3cm) node (s) {};
% % \shade[ball color=lightgray] ({2*5/3},0) arc (-90:25:.3cm) to [out=115,in=0] ({2*5/3},.35) node (b) {} to [out=180,in=65] (s.center) arc (155:270:.3cm);
% \path ({2*5/3-.1},0) to [out=170,in=-110]++ (-.25,.3) node (s1) {};
% \shade[ball color=lightgray] ({2*5/3+.1},0) to [out=10,in=-70]++ (.25,.3) to [out=110,in=0] ({2*5/3},.35) node (b) {} to [out=180,in=70] (s1.center) to [out=250,in=170] ({2*5/3-.1},0);

% %vec
% \draw[->,red,thick,vec] (a.center) --++ (0,-1) node[black,right] {$\vec F$};
% \node at (a.center) [circle,fill=red,inner sep=0pt,minimum size=1mm,label=above:Б] {};

% \draw[->,red,thick,vec] (b.center) --++ (0,-1) node[black,right] {$\vec F$};
% \node at (b.center) [circle,fill=red,inner sep=0pt,minimum size=1mm,label=above:Ш] {};

% \draw (b) --++ (-3,0);


    \end{tikzpicture}
\caption{Тела после снятия деформирующих сил}
\label{fig:p39}
\end{figure}

Как видно из рис.\,\ref{fig:p39}, деформации тел можно разделить на два вида.
\begin{itemize}
    \item \textbf{Упругая} деформация исчезает после снятия деформирующего действия.
    \item \textbf{Неупругая} деформация сохраняется (возможно, частично) после снятия деформирующей нагрузки.  
\end{itemize}

Таким образом, в рассмотренном примере в шаре возникали упругие деформации, а в бруске~--- неупругие. 

\sbox{0}{%
    \begin{tikzpicture}[font=\small,            line cap=round,                             >={Stealth[inset=0pt,round,length=3mm, angle'=20]},vec/.style={>={Stealth[inset=0pt,round,length=3.5mm, angle'=20]}}]


\filldraw[lightgray,semitransparent] (-0.25,0) rectangle (2.25,-0.3);
\draw (-0.25,0) -- (2.25,0);


\begin{scope}[xshift=.7cm]
    
\draw (0,0)
foreach \y in {.15,.45,...,3}
{ -- (.6,{\y}) -- (0,{\y+.15}) }
-- (.6,3)
;

\shade[ball color=lightgray] (.3,3.3) node (c) {} circle (.3);

%vec
\draw[->,red,thick,vec] (c.center) --++ (0,1) node[black,right] {$\vec F_\text{упр}$};
\node at (c.center) [circle,fill=red,inner sep=0pt,minimum size=1mm,label=above:] {};


% \draw[decoration = {zigzag,segment length = 2mm, amplitude = 2mm}, decorate] (1,0) -- ++(3,0);
% \draw (1,0) rectangle (1.2,.2);
\end{scope}


    \end{tikzpicture}%
}

{%
\setlength\intextsep{0pt}
\begin{wrapfigure}[17]{r}{\wd0}
    \centering
    \usebox{0}%
    \caption{Шар на пружине}
    \label{fig:p40}
\end{wrapfigure}


\textbf{Сила упругости} ($\vec F_\text{упр}$ [Н])~--- это сила, с которой упруго деформированное тело действует на соприкасающееся с ним тело, вызывающее деформацию. 

Пусть на пружину положили массивный шар (рис.\,\ref{fig:p40}). Под действием веса шара пружина деформируется (верхняя часть пружины смещается вниз), и возникает направленная вверх сила~$\vec F_\text{упр}$, удерживающая шар (остальные силы не показаны).



% \begin{figure}[H]
%     \centering

%     \begin{tikzpicture}[font=\small,            line cap=round,                             >={Stealth[inset=0pt,round,length=3mm, angle'=20]},vec/.style={>={Stealth[inset=0pt,round,length=3.5mm, angle'=20]}}]


% \filldraw[lightgray,semitransparent] (0,0) rectangle (5,-0.3);
% \draw (0,0) -- (5,0);

% \draw (.6,0) -- (0,0)
% foreach \y in {.15,.45,...,3}
% { -- (.6,{\y}) -- (0,{\y+.15}) }
% -- (.6,3)
% ;

% \shade[ball color=lightgray] (.3,3.3) circle (.3);


% % \draw[decoration = {zigzag,segment length = 2mm, amplitude = 2mm}, decorate] (1,0) -- ++(3,0);
% % \draw (1,0) rectangle (1.2,.2);

%     \end{tikzpicture}
% \caption{Тела после снятия деформирующих сил}
% \label{fig:p40}
% \end{figure}


Сила упругости:
\begin{itemize}
    \item направлена в
    сторону, \emph{противоположную смещению частиц} тела в процессе деформации;
    \item действует также \emph{между соседними слоями} деформированного тела и приложена к каждому слою.
\end{itemize}

}

\begin{law}
    \textbf{Закон Гука.} Сила упругости прямо пропорциональна величине деформации. Так, пружина, сжатая или растянутая на величину~$ x$, действует с силой
    \begin{equation}
    F_\text{упр}=kx,
    \end{equation}
    где~$k$~--- жёсткость пружины.
\end{law}

Стоит отметить, что сила реакции и вес представляют собой силы упругости, которые служат проявлением \emph{электромагнитного взаимодействия} тел.


 
\end{document}